\section{EXTRAÇÃO E APLICAÇÃO DOS NOVE DOCUMENTOS}
Foram reexaminados oito PDFs ABNT do repositório privado e recuperada a edição pública do IBGE correspondente à identificação de 1993 do nono documento. A cópia IBGE pública contém 60 páginas e não foi comparada binariamente ao arquivo privado. O registro de leitura conserva nomes, quantidade de páginas e hashes, sem redistribuir as normas.

Os documentos têm objetos diferentes: apresentação de trabalhos, referências, citações, resumos, numeração de seções, sumário, índice, lombada e tabelas. Sua combinação não forma automaticamente uma norma de artigo. A NBR 6022:2018, ausente da coleção privada, foi consultada adicionalmente no PDF hospedado pela UFRB para definir a estrutura específica do modelo demonstrativo (ABNT, 2018a).

\subsection{NBR 14724:2024 -- apresentação de trabalhos acadêmicos}
A estrutura distingue parte externa, com capa e lombada, e parte interna, com elementos pré-textuais, textuais e pós-textuais. As regras gerais tratam de formato, espaçamento, paginação, numeração, citações, siglas, equações, ilustrações e tabelas. Neste dossiê foram adotados capa, identificação, sumário, corpo organizado e referências, com margens de 3 cm superior/esquerda, 2 cm inferior/direita e entrelinhas de 1,5. Elementos de avaliação de grau, como folha de aprovação e ficha catalográfica, não foram inventados (ABNT, 2024, seções 4 e 5).

\subsection{NBR 6028:2021 -- resumo, resenha e recensão}
O documento organiza regras gerais e disposições para resumo, resenha e recensão. O modelo utiliza resumo informativo em parágrafo único, com objetivo, método, achados e limites. A extensão adotada está no intervalo recomendado para artigo, de 100 a 250 palavras. As palavras-chave são separadas por ponto e vírgula e encerradas por ponto (ABNT, 2021, seção 4.1).

\subsection{NBR 6023:2018, corrigida em 2020 -- referências}
A estrutura compreende elementos essenciais e complementares, localização, regras de apresentação, modelos por tipo documental, transcrição dos elementos e ordenação. A bibliografia dos documentos usa autores externos, dados de identificação e links verificáveis, com alinhamento à esquerda, espaço simples e separação entre entradas. O próprio corpus não é utilizado como fundamento bibliográfico do estudo. A coleção contém a edição histórica de 2018 corrigida em 2020; a UFES identifica atualização em 2025, cuja conferência integral permanece pendente (ABNT, 2018b; UFES, 2026).

\subsection{NBR 10520:2023 -- citações}
O documento trata de localização, regras gerais, sistemas de chamada, citações diretas e indiretas e notas. Foi adotado o sistema autor-data, com sobrenomes em grafia normal nas chamadas. As paráfrases são diferenciadas de resultados próprios. O modelo evita citações diretas extensas; se utilizadas futuramente, devem ter localização e apresentação próprias, sem substituir leitura da fonte por informação de catálogo (ABNT, 2023, seções 5 a 8).

\subsection{NBR 6024:2012 -- numeração progressiva}
A estrutura contém regras para seções, alíneas, subalíneas e indicativos. Os títulos recebem algarismos arábicos e hierarquia consistente, com espaço entre número e título. Resumos, sumário e referências não recebem indicativo numérico. O limite de subdivisão não exige criar níveis vazios: o modelo emprega somente as divisões necessárias ao argumento (ABNT, 2012a, seção 4).

\subsection{NBR 6027:2012 -- sumário}
O documento distingue localização, estrutura e regras de apresentação. O sumário do dossiê reproduz os títulos e a ordem do corpo do texto com seus localizadores de página; os elementos pré-textuais não são listados. O artigo demonstrativo não possui sumário, pois sua estrutura específica foi definida pela NBR 6022, e não pelo formato de monografia (ABNT, 2012b; ABNT, 2018a).

\subsection{NBR 6034:2004 -- índice}
A norma trata de definição, classificação, localização e apresentação do índice. Índice é instrumento de recuperação por entradas e localizadores, distinto de sumário. Não foi incluído índice remissivo nos dois documentos, porque não foi identificado um conjunto de entradas que o justificasse; o PDF permanece pesquisável (ABNT, 2004a, seções 3 a 6).

\subsection{NBR 12225:2004 -- lombada}
O arquivo denominado \emph{NBR-2012225-LOMBADA.pdf} contém internamente a NBR 12225:2004. Sua estrutura descreve lombada e título de margem, com autoria, título e elementos de identificação. A aplicação é própria de suporte encadernado; não foi gerada lombada para os PDFs. A edição antiga é preservada como identificação do acervo, sem ser declarada edição atual (ABNT, 2004b, seções 3 e 4).

\subsection{IBGE, 1993 -- apresentação tabular}
A edição lida organiza conceitos de tabela, espaços e elementos; elaboração; representação de tempo e classes de frequência; arredondamento; diagramação; recomendações; e exemplos. A informação central de uma tabela é numérica. Por isso, as comparações textuais deste estudo são apresentadas como quadros, e não como tabelas estatísticas. Uma futura tabela quantitativa deverá informar título, cabeçalho, unidades, fonte e notas, além de respeitar os critérios de moldura e continuidade (IBGE, 1993, seções 3, 4 e 8).

\subsection{Complemento indispensável: NBR 6022:2018}
A norma consultada define a sequência de identificação, autoria, resumos, elementos editoriais, introdução, desenvolvimento, considerações finais e referências. Seu item 6.1 recomenda fonte 12 e espaçamento simples para o artigo, deixando o projeto gráfico ao editor. Essa foi a correção principal em relação à versão anterior, que aplicava entrelinhas de trabalho acadêmico ao artigo. O modelo usa fonte da família Times, margens 3/2 cm como decisão gráfica e seções em 12 pontos; essas escolhas de família e margens não são apresentadas como imposições específicas da NBR 6022 (ABNT, 2018a, seções 5 e 6).

As datas editoriais são identificadas como não realizadas ou não aplicáveis, pois o artigo não foi submetido. Os contatos dos autores permanecem não informados. Essa lacuna não foi preenchida com dados inventados, e o documento é apresentado como modelo demonstrativo com análise piloto, sem alegação de conformidade integral com uma submissão final (ABNT, 2018a, seção 5.1).

\section{MATRIZ DE REQUISITOS E VERIFICAÇÃO}
A verificação relaciona documento, localizador, decisão de implementação e elemento visível no PDF. O Quadro 2 reúne as decisões principais; a conferência gráfica não substitui a revisão do conteúdo científico.

\begin{singlespace}
\begin{longtable}{|p{0.22\textwidth}|p{0.20\textwidth}|p{0.47\textwidth}|}
\caption{Requisitos extraídos e aplicação nos dois documentos.}\\\hline
Fonte & Localização & Aplicação \\\hline\endfirsthead
\multicolumn{3}{r}{\footnotesize Continuação do Quadro 2}\\\hline
Fonte & Localização & Aplicação \\\hline\endhead
6022:2018 & 5 e 6.1 & Artigo sem capa/sumário; pré-textuais, corpo e referências; fonte 12 e espaço simples.\\\hline
14724:2024 & 4 e 5.1--5.3 & Dossiê com capa, identificação, sumário; margens 3/2 cm, espaço 1,5 e paginação.\\\hline
6028:2021 & 4.1 & Resumo em parágrafo único e palavras-chave padronizadas.\\\hline
6023:2018/2020 & 6 e 9 & Referências externas, alinhadas à esquerda e ordenadas; revisão perante edição 2025 pendente.\\\hline
10520:2023 & 6.1 & Chamadas autor-data com grafia normal; achados próprios descritos sem autocitação.\\\hline
6024:2012 & 4.1 & Seções numeradas; elementos sem indicativo separados do corpo numerado.\\\hline
6027:2012 & 5 e 6 & Sumário somente no dossiê, coerente com títulos e páginas.\\\hline
6034:2004 & 4--6 & Índice remissivo não incluído; não confundido com sumário.\\\hline
12225:2004 & 3 e 4 & Lombada não aplicável ao suporte PDF entregue.\\\hline
IBGE, 1993 & 3, 4 e 8 & Comparações qualitativas como quadros; nenhuma estatística simulada.\\\hline
\end{longtable}
\noindent\footnotesize Fonte: síntese dos requisitos efetivamente consultados nos documentos normativos identificados nesta seção. As normas privadas não são reproduzidas como anexos.
\end{singlespace}
